% Standalone problem document, generated from the adjacent README.md.
% Compile with XeLaTeX. Mathematical content is maintained in Markdown.
\documentclass[11pt,a4paper]{article}
\usepackage[fontsize=10.5pt]{scrextend}
\usepackage[margin=22mm,headheight=15pt,headsep=9mm,footskip=11mm]{geometry}
\usepackage{fontspec}
\setmainfont{texgyrepagella}[Extension=.otf,UprightFont=*-regular,BoldFont=*-bold,ItalicFont=*-italic,BoldItalicFont=*-bolditalic]
\setsansfont{texgyreheros}[Extension=.otf,UprightFont=*-regular,BoldFont=*-bold,ItalicFont=*-italic,BoldItalicFont=*-bolditalic]
\setmonofont{texgyrecursor}[Extension=.otf,UprightFont=*-regular,BoldFont=*-bold,ItalicFont=*-italic,BoldItalicFont=*-bolditalic,Scale=MatchLowercase]
\usepackage{amsmath,amssymb}
\usepackage{unicode-math}
\setmathfont{texgyrepagella-math.otf}
\usepackage{xcolor,microtype,enumitem,needspace,fancyhdr}
\usepackage{bookmark}
\definecolor{ink}{HTML}{17344B}
\definecolor{accent}{HTML}{197D80}
\definecolor{muted}{HTML}{596773}
\hypersetup{colorlinks=true,linkcolor=accent,urlcolor=accent,pdftitle={IE-22 - Optimal
uniform row-deletion singular-value
constant},pdfauthor={Open Problems in Numerical Linear Algebra}}
\urlstyle{same}
\setlength{\parindent}{0pt}
\setlength{\parskip}{5pt plus 1pt minus 1pt}
\linespread{1.04}
\setlength{\emergencystretch}{3em}
\widowpenalty=10000
\clubpenalty=10000
\displaywidowpenalty=10000
\allowdisplaybreaks[1]
\setlist{nosep,leftmargin=1.5em}
\providecommand{\tightlist}{\setlength{\itemsep}{0pt}\setlength{\parskip}{0pt}}
\providecommand{\pandocbounded}[1]{#1}
\pagestyle{fancy}
\fancyhf{}
\fancyhead[L]{\sffamily\footnotesize\color{muted}OPEN PROBLEMS IN NUMERICAL LINEAR ALGEBRA}
\fancyhead[R]{\sffamily\bfseries\footnotesize\color{accent}IE-22}
\fancyfoot[L]{\parbox[b]{0.9\textwidth}{\sffamily\fontsize{8}{10}\selectfont\color{muted}Editorial ratings; a bounded search cannot certify that a problem remains unsolved.\\Literature check: 2026-09-10}}
\fancyfoot[R]{\sffamily\footnotesize\color{muted}\thepage}
\renewcommand{\headrulewidth}{0.3pt}
\renewcommand{\headrule}{\hbox to\headwidth{\color{accent}\leaders\hrule height \headrulewidth\hfill}}
\makeatletter
\renewcommand{\section}{\Needspace{5\baselineskip}\@startsection{section}{1}{0pt}{12pt plus 2pt minus 2pt}{4pt}{\sffamily\bfseries\large\color{ink}}}
\renewcommand{\subsection}{\Needspace{4\baselineskip}\@startsection{subsection}{2}{0pt}{9pt plus 1pt minus 1pt}{3pt}{\sffamily\bfseries\normalsize\color{ink}}}
\makeatother
\setcounter{secnumdepth}{0}
\begin{document}
{\sffamily\small\color{muted}Linear systems and elimination\par}
\vspace{3pt}
{\raggedright\hyphenpenalty=10000\exhyphenpenalty=10000\sffamily\bfseries\fontsize{21}{24}\selectfont\color{ink}Optimal
uniform row-deletion singular-value constant\par}
\vspace{7pt}
{\sffamily\small\textbf{Difficulty:} challenging\quad\textbf{Importance:} interesting
to the community\par}
{\sffamily\small\bfseries\color{accent}Status: Open\par}
\vspace{4pt}
{\color{accent}\hrule height 0.8pt}
\vspace{6pt}
\textbf{Rating rationale:} Challenging reflects an optimal uniform
extremum over matrices and deleted row sets; community impact is a sharp
robustness limit for row-sampling methods.

\subsection{Problem statement}\label{problem-statement}

Fix \(0<\theta<1\). For a real \(m\times n\) matrix \(A\) with every row
of Euclidean norm one, put

\[
s_\theta(A)=\min_{S\subseteq\{1,\ldots,m\},\ |S|=\lfloor\theta m\rfloor}
\min_{\|x\|_2=1}\|A_Sx\|_2,
\qquad
M_{m,n}(\theta)=\sup_{\|a_i\|_2=1}\sqrt{\frac nm}\,s_\theta(A).
\]

Here \(a_i\) is row \(i\) of \(A\), and \(A_S\) contains exactly the
rows indexed by \(S\).

Determine the sharp asymptotic upper constant for \(M_{m,n}(\theta)\) as
\(n\to\infty\) and \(m/n\to\infty\). In particular, is it

\[
c_\theta=\left(\frac1{\sqrt{2\pi}}\int_{-a_\theta}^{a_\theta}t^2e^{-t^2/2}\,dt\right)^{1/2},
\qquad
\mathbb P(|G|\le a_\theta)=\theta,\quad G\sim N(0,1)?
\]

Precisely, the proposed uniform inequality asks whether, for every
\(\varepsilon>0\), there exist integers \(N\) and \(R\) such that

\[
M_{m,n}(\theta)\le c_\theta+\varepsilon
\quad\text{whenever }n\ge N\text{ and }m/n\ge R,
\]

and whether no smaller constant has this property. The supremum and
quantifiers specify the source's uniform \(1+o(1)\) bound; the floor
convention handles row-count integrality.

\subsection{Connection to numerical linear
algebra}\label{connection-to-numerical-linear-algebra}

This asks how well any unit-row linear system can retain its smallest
singular value under worst-case row removal, a geometric constraint on
robust iterative solvers.

\newpage
\subsection{References}\label{references}

\begin{enumerate}
\def\labelenumi{\arabic{enumi}.}
\tightlist
\item
  S. Steinerberger, \emph{Quantile-based Random Kaczmarz for corrupted
  linear systems of equations}, Information and Inference \textbf{12}(1)
  (2023), 448--465, §2.3, the third question immediately after equation
  (8). \href{https://doi.org/10.1093/imaiai/iaab029}{Published paper}.
  \href{https://arxiv.org/abs/2107.05554}{Author preprint}, v1, §2.3,
  paragraph following \((\diamond)\).
\item
  E. Battaglia, J.-F. Cai, J. Chen, A. Ma, D. Needell, and T. Wu,
  \emph{Quantile Randomized Kaczmarz for Streaming Linear Systems with
  Massart Noise}, arXiv:2608.27968v1 (2026), §1.1 (background on the
  same restricted-singular-value heuristic and the different streaming
  model), §5. \href{https://arxiv.org/abs/2608.27968}{Preprint}.
\end{enumerate}

\subsection{Earlier status check ---
2026-09-08}\label{earlier-status-check-2026-09-08}

Checked the published question and latest arXiv records on 2026-09-08.
Targeted searches for ``Steinerberger subsingular'', ``small subsingular
values'', ``quantile Kaczmarz optimal constants'', ``unit rows smallest
singular value'', and 2025/2026 variants found no resolution of the
all-matrices extremum. Recent streaming Kaczmarz guarantees do not
supply this universal extremal inequality. The random-ensemble limit in
the same source is a related but distinct question: it specifies a
random construction, whereas this question optimizes over all unit-row
matrices. Openness remains subject to the limits of this search.

\subsection{Audit update --- 2026-09-10}\label{audit-update-2026-09-10}

Rechecked Steinerberger's
\href{https://arxiv.org/pdf/2107.05554}{primary formulation}, §2.3.
Searches for the optimal row-deletion constant and quantile-Kaczmarz
follow-ups found no proof of this extremal value. The new
\href{https://doi.org/10.1137/25M1785678}{Cai--Chen--Ma--Wu paper}, SIAM
J. Matrix Anal. Appl. 47 (2026), 802--823, proves a tight QRK
subsample-size result, not the universal matrix constant asked for here.
\end{document}
